% ---------------------------------------------------------------------------------------
%                  UOC - CSD - Spring 2025 - HY111 - Ask1 - csd5748 - :) 
% ---------------------------------------------------------------------------------------

\documentclass[fleqn]{article}
\setlength{\mathindent}{0pt}
\usepackage[utf8]{inputenc} 
\usepackage[english, greek]{babel}
\usepackage{subcaption}
\usepackage{hyperref}
\usepackage{xcolor}
\hypersetup
{
    hidelinks,  
    colorlinks=true,
    linkcolor=blue,
    urlcolor=blue 
}
\captionsetup[subfigure]{labelformat=empty}
\usepackage{amsmath}
\usepackage{amsfonts}
\usepackage{amssymb}
\usepackage{array}
\usepackage{graphicx}
\usepackage[a4paper, left=1.1cm, right=1.1cm, top=1.3cm, bottom=2cm]{geometry}
\usepackage[absolute, overlay]{textpos}
\usepackage{listings}
\usepackage{xcolor}
\setlength{\jot}{0pt}
\lstdefinestyle{mypython}
{
    language=Python,
    basicstyle=\ttfamily\small,
    keywordstyle=\color{blue},
    stringstyle=\color{red},
    commentstyle=\color{green!50!black},
    numbers=left,
    numberstyle=\tiny\color{gray},
    stepnumber=1,
    frame=single,
    backgroundcolor=\color{gray!10},
    breaklines=true
}
\begin{document}
\selectlanguage{greek}
% ---------------------------------------------------------------------------------------

\begin{textblock}{10}(-0.299,0.085)
    \boxed{\textbf{\LaTeX{}}}
\end{textblock}

\begin{center}
\setlength{\fboxsep}{10pt}
\noindent \boxed{\boxed{\textbf{\Large    ΗΥ-111: Απειροστικός Λογισμός ΙI:  Άσκηση 1\textsuperscript{η}  }}}
\end{center}

% ---------------------------------------------------------------------------------------

\vspace{0.999cm}
\begin{center}
\boxed{\underline{\textbf{· 1\textsuperscript{η} Άσκηση -}\textbf{ Ερώτημα (α):}}}
\end{center}
\vspace{0.321cm}

\[
    \textbf{Δεδομένα:} \hspace{0.321cm} \vec{v} = \vec{i} + \vec{j} + 2\vec{k} \hspace{0.321cm} \textbf{και} \hspace{0.321cm} \vec{u} = -\vec{i} - \vec{k} \hspace{0.321cm} (\mathbb{R}^3)  \hspace{0.321cm} 
\]

\[
    [i][ii] \hspace{0.321cm} \textbf{Ξέρουμε:} \hspace{0.321cm} ||\vec{u}|| = \sqrt{u_1^2 + u_2^2 + u_3^2} \hspace{0.321cm} \textbf{το μέτρο ή μήκος του διανύσματος} \hspace{0.321cm} \vec{u} = u_1\cdot\vec{i} + u_2 \cdot \vec{j} + u_3 \cdot \vec{k}
\]

\[
    \textbf{Οπότε είναι:} \hspace{0.321cm} ||\vec{v}|| = \sqrt{1^2 + 1^2 + 2^2} = \sqrt{6} \hspace{0.321cm} \textbf{και} \hspace{0.321cm} ||\vec{u}|| = \sqrt{(-1)^2 + 0^2 + (-1)^2} = \sqrt{2}
\]

\[
    [iii] \hspace{0.321cm} \textbf{Ξέρουμε:} \hspace{0.321cm} ||\vec{u}|| \cdot ||\vec{v}|| = |u|\cdot|v|\cdot cos(\theta) \hspace{0.321cm} \textbf{ή} \hspace{0.321cm} ||\vec{u}|| \cdot ||\vec{v}|| = u_1\cdot v_1 + u_2\cdot v_2 + u_3\cdot v_3
\]

\[
    \textbf{το εσωτερικό γινόμενο των διανυσμάτων} \hspace{0.321cm} \vec{u} = u_1\cdot\vec{i} + u_2\cdot\vec{j} + u_3\cdot\vec{k} \hspace{0.321cm} \textbf{και} \hspace{0.321cm} \vec{v} = v_1\cdot\vec{i} + v_2\cdot\vec{j} + v_3\cdot\vec{k}
\]

\[
    \textbf{Οπότε είναι:} \hspace{0.321cm} 
    v \cdot u = (1)(-1) + (1)(0) + (2)(-1) = -1 + 0 - 2 = -3 \hspace{0.321cm} 
\]

\[
    [iv] \hspace{0.321cm} \textbf{Ξέρουμε:} \hspace{0.321cm} \vec{u} \times \vec{v} = ( | \vec{u} |\cdot| \vec{v} |\cdot sin(\theta) ) \cdot \vec{n} 
    \hspace{0.321cm} \textbf{ή} \hspace{0.321cm}
\]

\[
    \vec{u} \times \vec{v} = 
    \begin{vmatrix}
    \vec{i} & \vec{j} & \vec{k} \\
    u_1 & u_2 & u_3 \\
    v_1 & v_2 & v_3
    \end{vmatrix} =
    \begin{vmatrix}
    u_2 & u_3 \\
    v_2 & v_3
    \end{vmatrix} \cdot \vec{i} -
    \begin{vmatrix}
    u_1 & u_3 \\
    v_1 & v_3
    \end{vmatrix} \cdot \vec{j} +
    \begin{vmatrix}
    u_1 & u_2 \\
    v_1 & v_2
    \end{vmatrix} \cdot \vec{k} =
    (u_2\cdot v_3 - u_3 \cdot v_2)\cdot \vec{i} -
    (u_1\cdot v_3 - u_3 \cdot v_1)\cdot \vec{j} +
    (u_1\cdot v_2 - u_2 \cdot v_1)\cdot \vec{k}
    \hspace{0.123456789cm}
\]

\[
    \textbf{το εξωτερικό γινόμενο των διανυσμάτων} \hspace{0.321cm} \vec{u} = u_1\cdot\vec{i} + u_2\cdot\vec{j} + u_3\cdot\vec{k} \hspace{0.321cm} \textbf{και} \hspace{0.321cm} \vec{v} = v_1\cdot\vec{i} + v_2\cdot\vec{j} + v_3\cdot\vec{k}
\]

\[
    \textbf{Οπότε είναι:} \hspace{0.321cm} 
    \vec{u} \times \vec{v} = 
    \begin{vmatrix}
    \vec{i} & \vec{j} & \vec{k} \\
    -1 & 0 & -1 \\
    1 & 1 & 2
    \end{vmatrix}
    = (0\cdot2 - (-1)\cdot1)\cdot\vec{i} - (-1\cdot2 - (-1)\cdot1)\cdot\vec{j} + (-1\cdot1 - 0\cdot1)\cdot\vec{k} = \vec{i} + \vec{j} - \vec{k}
\]

\[
    [v] \hspace{0.321cm} \textbf{Ξέρουμε:} \hspace{0.321cm} 
    \vec{v} \times \vec{u} = - (\vec{u} \times \vec{v}) = - (\vec{i} + \vec{j} - \vec{k})
    = -\vec{i} - \vec{j} + \vec{k}
\]

\[
    [vi] \hspace{0.321cm} \textbf{Από } [i][ii] \textbf{ έχουμε:} \hspace{0.321cm} 
    ||\vec{v} \times \vec{u}|| = \sqrt{(-1)^2 + (-1)^2 + (+1)^2} = \sqrt{3}
\]

\newpage

\[
    [vii] \hspace{0.321cm} \textbf{Ξέρουμε:} \hspace{0.321cm} 
    \cos(\theta) = \frac{\vec{u} \cdot \vec{v}}{||\vec{u}|| \cdot ||\vec{v}||}
\]

\[
    \textbf{Οπότε είναι:} \hspace{0.321cm} 
    \cos(\theta) = \frac{\vec{u} \cdot \vec{v}}{||\vec{u}|| \cdot ||\vec{v}||} =
    \frac{\vec{v} \cdot \vec{u}}{||\vec{u}|| \cdot ||\vec{v}||}
    = \frac{-3}{\sqrt{2} \cdot \sqrt{6}} = \frac{-3}{\sqrt{12}} = \frac{-3}{2\sqrt{3}} 
    = - \frac{\sqrt{3}}{2} \hspace{0.321cm} \textbf{και άρα} \hspace{0.321cm} 
    \hat \theta = \cos^{-1} \left( -\frac{\sqrt{3}}{2} \right) = \frac{5\pi}{6}
\]

\[
    [viii] \hspace{0.321cm} \textbf{Ξέρουμε:} \hspace{0.321cm} 
    proj_v(\vec{u}) = ( ||\vec{u}||\cdot \cos(theta) ) \cdot \frac{\vec{v}}{||\vec{v}||} = (\frac{\vec{u} \cdot \vec{v}}{||\vec{v}||} \cdot \frac{\vec{v}}{||\vec{v}||} ) = \left( \frac{\vec{u} \cdot \vec{v}}{||\vec{v}||^{2}} \right) \cdot \vec{v}
\]

\[
    \textbf{Οπότε είναι:} \hspace{0.321cm} 
    \left( \frac{\vec{u} \cdot \vec{v}}{||\vec{v}||^{2}} \right) \cdot \vec{v} =
    \frac{-3}{\sqrt{6}^2} \cdot (\vec{i} + \vec{j} + 2\vec{k}) = -\frac{1}{2} (\vec{i} + \vec{j} + 2\vec{k}) = -\frac{1}{2}\vec{i} - \frac{1}{2}\vec{j} - \vec{k}
\]

\vspace{1.555cm}
\begin{center}
\boxed{\underline{\textbf{· 1\textsuperscript{η} Άσκηση -}\textbf{ Ερώτημα (β):}}}
\end{center}
\vspace{0.5cm}


\setlength{\fboxrule}{1pt}
\begin{figure}[h]
  \hspace{-0.3cm}
  \centering
  \fbox{\includegraphics[width=0.99\textwidth]{geogebra1.png}}
\end{figure}

\[
    \textbf{Λοιπόν, σε αυτήν την εικόνα από το \textlatin{GeoGebra}, μπορούμε να δούμε γραφηκά όλα τα ζητούμενα του} 
\]
\[
    \textbf{ ερωτήματος (α), με την ίδια σειρά. Έτσι, έχουμε πως } ||\vec{v}|| = 2.449 \approx \sqrt{6} \textbf{ και } ||\vec{u}|| = 1.414 \approx \sqrt{2} \textbf{ σωστά !} \hspace{1.555cm}
\]
\[
    \textbf{Επίσης, } ||\vec{v} \times \vec{u}|| \overset{*}{=} |d| = 1.732 \approx \sqrt{3} \textbf{ και } \cos{({\vec{u}, \vec{v})}} = 150^o = \frac{5\pi}{6} \textbf{ και τα υπόλοιπα που δεν ανέφερα είναι οι συ-}
\]
\[
    \textbf{-ντελεστές των διανυσμάτων } \vec{i}, \vec{j}, \vec{k} \textbf{ αντοίστιχα. Τελικά, όλα τα αποτελέσματα επαληθεύτηκαν ..!} \hspace{1cm}
\]

% ---------------------------------------------------------------------------------------

\newpage

\vspace{1.555cm}
\begin{center}
\boxed{\underline{\textbf{· 2\textsuperscript{η} Άσκηση:}}}
\end{center}
\vspace{0.321cm}

\[
    \textbf{Δεδομένα:} \hspace{0.321cm} 
    \vec{u} = 2\vec{i} + 4\vec{j} - 5\vec{k} 
    \hspace{0.321cm} \textbf{και} \hspace{0.321cm}
    \vec{v} = -4\vec{i} - 8\vec{j} + \alpha\vec{k}
    \hspace{0.321cm} (\mathbb{R}^3)  \hspace{0.321cm} 
\]

\[
    \textbf{Ξέρουμε πως δύο διανύσματα είναι παράλληλα αν και μόνο αν το εξωτερικό τους γινόμενο}
\]

\[
    \textbf{είναι το μηδενικό διάνυσμα, ή} \hspace{0.321cm} \vec{u} \times \vec{v} = \vec{0}
\]

\[
    \textbf{Οπότε είναι:} \hspace{0.321cm} 
    \vec{u} \times \vec{v} = 
    \begin{vmatrix}
    \vec{i} & \vec{j} & \vec{k} \\
    2 & 4 & -5 \\
    -4 & -8 & \alpha
    \end{vmatrix}
    = 
    \begin{vmatrix} 
    4 & -5 \\ 
    -8 & \alpha 
    \end{vmatrix} \vec{i} 
    - 
    \begin{vmatrix} 
    2 & -5 \\ 
    -4 & \alpha 
    \end{vmatrix} \vec{j} 
    + 
    \begin{vmatrix} 
    2 & 4 \\ 
    -4 & -8 
    \end{vmatrix}\vec{k} \hspace{0.321cm} \implies
\]

\[
    \Rightarrow (4\alpha - (-5) \cdot (-8)) \cdot \vec{i} - (2\alpha - (-5) \cdot (-4)) \cdot \vec{j} + (2 \cdot (-8) - 4 \cdot (-4)) \cdot \vec{k} = (4\alpha - 40) \cdot \vec{i} - (2\alpha - 20) \cdot \vec{j} + (0) \cdot \vec{k}
\]

\[
    \textbf{Και αφού θέλουμε να μηδενιστεί, πρέπει:}
\]

\[
    4\alpha - 40 = 0 \hspace{0.321cm} \textbf{και} \hspace{0.321cm} 2\alpha - 20 = 0
    \hspace{0.321cm} \textbf{και} \hspace{0.321cm} (0) = 0
    \hspace{0.321cm} \implies \hspace{0.321cm}
    \alpha = 10 \hspace{0.321cm} \textbf{και} \hspace{0.321cm} a = 10
    \hspace{0.321cm} \textbf{και} \hspace{0.321cm} 0 = 0 \textbf{ (ισχύει) }
\]

\[
    \textbf{Τελικά,} \hspace{0.321cm} \boxed{ \alpha = 10 }
\]

% ---------------------------------------------------------------------------------------

\vspace{0.555cm}
\begin{center}
\boxed{\underline{\textbf{· 3\textsuperscript{η} Άσκηση:}}}
\end{center}
\vspace{0.321cm}

\[
    \textbf{Δεδομένα:} \hspace{0.321cm} 
    \mathbb{R}^3, \quad \vec{a}, \quad \vec{b}, \quad ||\vec{a}|| = ||\vec{b}|| = 1, \quad
    { \cos ( \vec{a}, \vec{b}}) = 60^o = \frac{1}{2}, \quad \vec{x} \implies \vec{x} \parallel (\vec{b}-\vec{a}),
    \quad \vec{a} \perp (b-2\vec{x}).
\]

\[
    \textbf{Αρχικά:} \hspace{0.321cm} 
    \vec{a} \cdot \vec{a} = (\vec{a})^2 = (||\vec{a}||)^2 = 1^2 = 1
    \hspace{0.321cm} \textbf{και} \hspace{0.321cm}
    \vec{a} \cdot \vec{b} = ||\vec{a}|| \cdot ||\vec{b}|| \cdot \cos({\vec{a}, \vec{b}})
    = 1 \cdot 1 \cdot \frac{1}{2} = 0.5
\]

\[
    \textbf{Επίσης, αφού:} \hspace{0.321cm} 
    \vec{x} \parallel (\vec{b}-\vec{a})
    \hspace{0.321cm} \textbf{τότε} \hspace{0.321cm}
    \vec{x} = k \cdot (\vec{b} - \vec{a}) \hspace{0.321cm} (1)
    \hspace{0.321cm} \textbf{όπου \textlatin{k} ένας πραγματικός αριθμός.} \hspace{0.321cm}
\]

\[
    \textbf{Επιπλέον, αφού:} \hspace{0.321cm} 
    \vec{a} \perp (b-2\vec{x})
    \hspace{0.321cm} \textbf{τότε} \hspace{0.321cm}
    \vec{a} \cdot (\vec{b}-2\vec{x}) = 0 \hspace{0.321cm} (2)
\]

\[
    \textbf{Έτσι, αντικαθιστώντας την (1) στην (2):} \hspace{0.321cm} 
    \vec{a} \cdot (\vec{b}-2 \cdot (k \cdot (\vec{b} - \vec{a}))) = 0 \hspace{0.123cm} \Rightarrow
    \hspace{0.123cm} \vec{a} \cdot \vec{b} - 2k \cdot (\vec{a} \cdot \vec{b} - \vec{a} \cdot \vec{a} ) = 0
\]

\[
    \implies \frac{1}{2} - 2k \cdot (\frac{1}{2} - 1) = 0 \Rightarrow k = -\frac{1}{2}
\]

\[
    \textbf{Τελικά, ξανά στην σχέση (1)} \hspace{0.321cm} \vec{x} = -\frac{1}{2} \cdot (\vec{b} - \vec{a}) = \boxed{\frac{1}{2}\vec{a} - \frac{1}{2}\vec{b} = \vec{x}}
\]

% ---------------------------------------------------------------------------------------

\newpage

\vspace{0.5cm}
\begin{center}
\boxed{\underline{\textbf{· 4\textsuperscript{η} Άσκηση -}\textbf{ Ερώτημα (α):}}}
\end{center}
\vspace{0.1cm}

\[
    \textbf{Δεδομένα:} \hspace{0.321cm} 
    \epsilon_1: 12x + 5y = 1 \hspace{0.321cm}\textbf{και}\hspace{0.321cm} \epsilon_2 : 2x - 2y = 3.
\]

\[
    \textbf{Ξέρουμε, πως το διάνυσμα} \hspace{0.321cm} 
    \vec{u} = a \vec{i} + b \vec{j} 
    \hspace{0.321cm} \textbf{είναι κάθετο στην ευθεία} \hspace{0.321cm}
    \epsilon : ax + by = c
\]

\[
    \textbf{(αφού ισχύει πως} \hspace{0.321cm} \lambda_{\vec{u}} \cdot \lambda_{\epsilon} =
    -1 \Rightarrow \frac{b}{a} \cdot (-\frac{b}{a}) = -1 \Rightarrow -1 = -1 \textbf{)}
\]

\[
    \textbf{Οπότε είναι:} \hspace{0.321cm} 
    \vec{u}_1 = 12\vec{i} + 5\vec{j} \hspace{0.321cm} \textbf{και} \hspace{0.321cm} 
    \vec{u}_2 = 2\vec{i} - 2\vec{j} \hspace{0.321cm} \textbf{(τα δύο κάθ. στις αντιστ. ευθ. διαν.)}
\]

\[
    \textbf{Έτσι, τώρα από τον τύπο:} \hspace{0.321cm} 
    \cos(\theta) = \frac{\vec{u_1} \cdot \vec{u_2}}{||\vec{u_1}|| \cdot ||\vec{u_2}||} 
    \hspace{0.321cm} \overset{*}{\implies} \hspace{0.321cm} \cos(\theta) = \frac{14}{13 \cdot 2\sqrt{2}} = \frac{7}{13\sqrt{2}}
\]

\[
    \bold{*} \hspace{0.321cm} \vec{u_1} \cdot \vec{u_2} = (12)(2) + (5)(-2) = 14
\]

\[
    \bold{*} \hspace{0.321cm} ||\vec{u_1}|| = \sqrt{12^2 + 5^2} = \sqrt{169} = 13 
    \hspace{0.321cm} \textbf{και} \hspace{0.321cm}
    ||\vec{u_2}|| = \sqrt{2^2 + (-2)^2} = \sqrt{8} = 2\sqrt{2}
\]

\[
    \textbf{Τελικά:} \hspace{0.321cm}
    \theta = \cos^{-1} (\frac{7}{13\sqrt{2}}) \hspace{0.321cm} \implies \hspace{0.321cm}
    \theta \approx 67.5^o \Rightarrow \theta \approx 1.18 \hspace{0.12345cm} rad
    \hspace{0.321cm} \textbf{( ή \guillemotleft \hspace{0.1cm}a \guillemotright \hspace{0.1cm} παρακάτω στο ερώτημα (β) )}
\]

\vspace{0.5cm}
\begin{center}
\boxed{\underline{\textbf{· 4\textsuperscript{η} Άσκηση -}\textbf{ Ερώτημα (β):}}}
\end{center}
\vspace{0.1cm}

\setlength{\fboxrule}{1pt}
\begin{figure}[h]
  \hspace{-0.3cm}
  \centering
  \fbox{\includegraphics[width=0.99\textwidth]{geogebra2.jpg}}
\end{figure}

\vspace{0.1cm}

% ---------------------------------------------------------------------------------------

\newpage

\textbf{ }
\vspace{0.5cm}
\begin{center}
\boxed{\underline{\textbf{· 5\textsuperscript{η} Άσκηση -}\textbf{ Ερώτημα (α):}}}
\end{center}
\vspace{0.25cm}

\[
    \textbf{Δεδομένα:} \hspace{0.321cm} 
    Ο(0, 0, 0) \hspace{0.321cm} \textbf{και} \hspace{0.321cm} 
    \textbf{επ.} \hspace{0.1cm} \perp \hspace{0.1cm} 2x+3y+z=12
\]

\[
    \textbf{Ξέρουμε, πως η γενική μορφή ενός επιπέδου στο τρισδιάστατο χώρο είναι:} \hspace{0.321cm} 
\]

\[
    Ax+By+Cz+D=0
\]

\[
    \textbf{Όμως, το επίπεδο που ψάχνουμε περνάει από την αρχή των αξόνων, οπότε:} \hspace{0.2cm} D=0
\]

\[
    \textbf{Επίσης, το διάνυσμα του επιπέδου που μας δίνεται είναι:}
\]

\[
    \vec{u_1} = Α\vec{i} + B\vec{j} + C\vec{k} \hspace{0.321cm} \implies \hspace{0.321cm}
    \vec{u_1} = 2\vec{i} + 3\vec{j} + 1\vec{k}
\]

\[
    \textbf{Ετσι, αν βρούμε ένα διάνυσμα που να είναι κάθετο σε αυτό, θα βρούμε και το ζητούμενο.}
\]

\[
    \vec{u_2} \perp \vec{u_1} \Rightarrow \vec{u_2} \cdot \vec{u_1} = 0 \Rightarrow
    2A + 3B + 1C = 0 \Rightarrow \infty \textbf{ λύσεις ...}
\]

\[
    \textbf{Οπότε, μία από τις άπειρες λύσεις είναι και το διάνυσμα:} \hspace{0.321cm}
    \vec{u_2} = \vec{i} + \vec{j} -5\vec{k} 
\]

\[
    \textbf{Τελικά, ένα επίπεδο θα είνια και το:} \hspace{0.321cm}
    \boxed{x + y -5z = 0} \hspace{0.321cm} (A=1, B=1, C=-5, D=0)
\]





\vspace{0.75cm}
\begin{center}
\boxed{\underline{\textbf{· 5\textsuperscript{η} Άσκηση -}\textbf{ Ερώτημα (β):}}}
\end{center}
\vspace{0.25cm}

\[
    \textbf{Δεδομένα:} \hspace{0.321cm} 
    \textbf{επ1:} \hspace{0.2cm} 2x+3y+z=12
    \hspace{0.321cm} \textbf{και} \hspace{0.321cm} 
    \textbf{επ2:} \hspace{0.2cm} x + y -5z = 0
\]

\[
    \textbf{Ξέρουμε, πως η (διανυσματική, όχι παραμετρική) εξίσωση ευθείας (τομής επιπέδων) είναι:} \hspace{0.321cm} 
\]

\[
    \cdot \hspace{0.321cm} \vec{r}(t)= \vec{r_0} + t\vec{v}
\]

\[
    \textbf{Όπου} \hspace{0.25cm} r_0 \hspace{0.25cm} \textbf{το διάνυσμα θέσης (δηλ. από αρχή αξόνων έως σημείο} \hspace{0.25cm} P_0(x_0,y_0,z_0) \textbf{) ...}
\]

\vspace{1cm}
\newpage

\[
    \textbf{Και} \hspace{0.25cm} \vec{v} \hspace{0.25cm} \textbf{ένα διάνυσμα που είναι παράλληλο στην ευθεία} \hspace{0.25cm} \vec{L} \hspace{0.321cm} \textbf{(και} \hspace{0.1cm} -\infty < t < +\infty
 \hspace{0.123cm} \textbf{) ...}
\]

\[
\textbf{Τώρα, ένα} \hspace{0.321cm} \vec{r_0} \hspace{0.321cm} \textbf{θα είναι το:}
\hspace{0.321cm}
\left\{
\begin{array}{l}
\cdot \hspace{0.321cm} 2x+3y+z=12 \\
\cdot \hspace{0.321cm} x + y -5z = 0
\end{array}
\right.
\hspace{0.321cm}
\implies
\hspace{0.321cm}
\left\{
\begin{array}{l}
\cdot \hspace{0.321cm} 2x+3y+z-12=0 \\
\cdot \hspace{0.321cm} x + y -5z = 0
\end{array}
\right.
\]

\[
    \textbf{Όμως, το συγκεκριμενο σύστημα έχει άπειρες λύσεις. Συγκεκριμένα (γενικά):}
\]

\[
\left\{
\begin{array}{l}
\cdot \hspace{0.321cm} x=16z-12 \\
\cdot \hspace{0.321cm} y=12-11z \\
\cdot \hspace{0.321cm} z = z
\end{array}
\right.
\hspace{0.321cm} \implies \hspace{0.321cm}
\textbf{π.χ.} \hspace{0.321cm} x=-12, \hspace{0.1cm} y=12, \hspace{0.1cm} z=0
\hspace{0.321cm} \implies \hspace{0.321cm} \underline{\vec{r_0} = -12\vec{i} + 12\vec{j}}
\]

\[
    \textbf{Επιπλέον, ένα} \hspace{0.321cm} \vec{v} \hspace{0.321cm} \textbf{είναι το}
    \hspace{0.321cm} \vec{v} = \vec{u_1} \times \vec{u_2} \hspace{0.321cm} 
    \textbf{(} \vec{v} \perp \vec{u_1} \hspace{0.25cm}\textbf{όμως}\hspace{0.25cm} \vec{u_1} \perp \textbf{ επ1 } \perp \textbf{ επ2 } \perp \vec{L} \hspace{0.25cm}\textbf{άρα}\hspace{0.25cm}
    \vec{v} \parallel \vec{L} \textbf{)}
\]

\[
    \textbf{Έτσι:} \hspace{0.321cm}
    \vec{v} \hspace{0.321cm} = \hspace{0.321cm}
    \begin{vmatrix}
    \vec{i} & \vec{j} & \vec{k} \\
    2 & 3 & 1 \\
    1 & 1 & -5
    \end{vmatrix}
    \hspace{0.321cm} = \hspace{0.321cm} \cdot \cdot \cdot \hspace{0.321cm} \implies \hspace{0.2cm} \underline{\vec{v} \hspace{0.1cm} = \hspace{0.1cm} -16\vec{i} + 11\vec{j} - 1\vec{k}}
\]

\[
    \textbf{Τελικά:} \hspace{0.321cm}
    \vec{r}(t) = \hspace{0.1cm}<-12, 12, 0>\hspace{0.1cm} + \hspace{0.1cm}t \hspace{0.1cm}\cdot <-16 + 11 - 1>
    \hspace{0.111cm} \implies \hspace{0.111cm}
    \boxed{\vec{r}(t) = (-12-16t)\vec{i} + (12+11t)\vec{j} +(-t)\vec{k}}
\]





\vspace{0.99cm}
\begin{center}
\boxed{\underline{\textbf{· 5\textsuperscript{η} Άσκηση -}\textbf{ Ερώτημα (γ):}}}
\end{center}
\vspace{0.25cm}

\setlength{\fboxrule}{1pt}
\begin{figure}[h]
  \hspace{-0.3cm}
  \centering
  \fbox{\includegraphics[width=0.99\textwidth]{geogebra3.jpg}}
\end{figure}

% ---------------------------------------------------------------------------------------

\newpage

\textbf{ }
\vspace{0.5cm}
\begin{center}
\boxed{\underline{\textbf{· 6\textsuperscript{η} Άσκηση -}\textbf{ Ερώτημα (α):}}}
\end{center}
\vspace{0.25cm}

\[
    \textbf{Δεδομένα:} \hspace{0.321cm} 
    \epsilon : \hspace{0.2cm} x = 3-2t,\hspace{0.1cm} y = 1+5t,\hspace{0.1cm} z = -2-3t
    \hspace{0.321cm} \textbf{και} \hspace{0.321cm} 
    \epsilon \pi. \hspace{0.2cm} 2x+y-z=17 .
\]

\[
    \textbf{Ξέρουμε, πως οι παραμετρικές εξισώσεις μίας ευθείας είναι:} \hspace{0.321cm} 
\]

\[
    \cdot \hspace{0.321cm} x = x_0 + t v_1, \hspace{0.321cm} y = y_0 + t v_2, \hspace{0.321cm} z = z_0 + t v_3 \hspace{0.321cm}  (-\infty < t < +\infty)
\]

\[
    \textbf{Άρα, αν} \hspace{0.321cm} \vec{v} = v_1\vec{i} + v_2\vec{j} + v_3\vec{k}
    \hspace{0.321cm} \textbf{θα έχουμε} \hspace{0.321cm}
    \vec{v} = -2\vec{i} + 5\vec{j} -3\vec{k}
\]

\[
    \textbf{Επίσης, το διάνυσμα του επιπέδου είναι:} \hspace{0.321cm} 
    \vec{u} = 2\vec{i} + 1\vec{j} + -1\vec{k}
\]

\[
    \textbf{Έτσι, πρέπει:} \hspace{0.321cm} 
    \vec{v} \cdot \vec{u} = 0 \Rightarrow
    (-2,5,-3) \cdot (2,1,-1)= 0 \Rightarrow (-2)(2)+(5)(1)+(-3)(-1)=0 \Rightarrow 4=0
\]

\[
    \textbf{Άτοπο, οπότε αφού} \hspace{0.25cm} \vec{v} \not\perp \vec{u} 
    \hspace{0.25cm} \textbf{τότε} \hspace{0.25cm} 
    \boxed{\epsilon \not\parallel \epsilon \pi.}
\]

\[
    \textbf{Ακολουθεί και μια ακόμη (άτυπη) εικονική αιτιολόγηση από το} \hspace{0.2cm}
    \href{https://www.geogebra.org}{\textcolor{blue}{Geogebra}} \hspace{0.2cm} ...
\]

\vspace{0.75cm}
\setlength{\fboxrule}{1pt}
\begin{figure}[h]
  \hspace{-0.3cm}
  \centering
  \fbox{\includegraphics[width=0.99\textwidth]{geogebra4.jpg}}
\end{figure}


% ---------------------------------------------------------------------------------------

\newpage

\textbf{ }
\vspace{0.1cm}
\begin{center}
\boxed{\underline{\textbf{· 7\textsuperscript{η} Άσκηση -}\textbf{ Ερώτημα (α και β):}}}
\end{center}
\vspace{0.35cm}

\[
    \textbf{Δεδομένα:} \hspace{0.321cm} 
    \epsilon \pi : \hspace{0.2cm} x+2y+6z = 6
    \hspace{0.321cm} \textbf{και} \hspace{0.321cm} 
    \epsilon : \hspace{0.2cm} x = 2+t,\hspace{0.2cm} y = 1+t, \hspace{0.2cm}z = - \frac{1}{2} - \frac{1}{2} t .
\]

\[
    \textbf{Ξέρουμε, η απόσταση ενός σημείου από ένα επίπεδο στο χώρο:} \hspace{0.321cm} 
    d \hspace{0.1cm} = \hspace{0.1cm} | \hspace{0.1cm} \frac{\vec{PS} \cdot \vec{n}}{||\vec{n}||} \hspace{0.1cm} |
\]

\[
    \textbf{Η αλλιώς:} \hspace{0.321cm} 
    d \hspace{0.1cm} = \hspace{0.1cm} | \hspace{0.1cm} \frac{P \cdot \epsilon \pi .}{||\vec{n}||} \hspace{0.1cm} | \hspace{0.321cm} \textbf{όπου} \hspace{0.321cm}
    P=(x_0,y_0,z_0) \hspace{0.321cm} \textbf{τυχαίο σημείο.}
\]

\[
    \textbf{Τελικά:} \hspace{0.321cm} 
    d \hspace{0.1cm} = \hspace{0.1cm} | \hspace{0.1cm} 
    \frac{Ax_0 + By_0 + Cz_0 + D}{\sqrt{A^2 + B^2 + C^2 }} \hspace{0.1cm} | \hspace{0.321cm}
\]

\[
    \textbf{Και έχουμε:} \hspace{0.321cm} 
    \epsilon \pi : \hspace{0.2cm} x+2y+6z = 6 \hspace{0.2cm} \implies \hspace{0.2cm}
    \underline{ A=1,\hspace{0.2cm}B=2,\hspace{0.2cm}C=6,\hspace{0.2cm}D=-6 }
\]

\[
    \textbf{Επιπλέον, ένα τυχαίο σημείο της ευθείας αυτής είναι:} \hspace{0.321cm} 
\]

\[
    \textbf{Έστω} \hspace{0.321cm} t \hspace{0.1cm} = \hspace{0.1cm} 0 \hspace{0.321cm} \textbf{τότε} \hspace{0.321cm} x \hspace{0.1cm} = \hspace{0.1cm} 2, \hspace{0.1cm} y \hspace{0.1cm} = \hspace{0.1cm} 1, \hspace{0.1cm} z \hspace{0.1cm} = \hspace{0.1cm} -\frac{1}{2} \hspace{0.321cm} \textbf{και έτσι} \hspace{0.321cm} \underline{P \hspace{0.1cm} = \hspace{0.1cm} (2, \hspace{0.1cm} 1, \hspace{0.1cm} -1/2)}
\]

\[
    \textbf{Άρα:} \hspace{0.4321cm}
    d = | \frac{1(2) + 2(1) + 6(-\frac{1}{2}) + (-6)}{\sqrt{1^2+ 2^2 + 6^2}} |
    = | \frac{-5}{\sqrt{41}} | = \frac{5}{\sqrt{41}}
    \hspace{0.321cm} \implies \hspace{0.321cm} \boxed{ d = \frac{5\sqrt{41}}{41} \approx 0.780 }
\]


\vspace{0.55cm}
\setlength{\fboxrule}{1pt}
\begin{figure}[h]
  \hspace{-0.3cm}
  \centering
  \fbox{\includegraphics[width=0.99\textwidth]{geogebra5.jpg}}
\end{figure}

\vspace{0.35cm}


% ---------------------------------------------------------------------------------------

\newpage

\textbf{ }
\vspace{0.1cm}
\begin{center}
\boxed{\underline{\textbf{· 8\textsuperscript{η} Άσκηση -}\textbf{ Ερώτημα (όλα):}}}
\end{center}
\vspace{0.3cm}

\[
    \alpha) \hspace{0.321cm} proj_v(\vec{u}) = proj_u(\vec{v}) \Rightarrow 
    \frac{\vec{u} \cdot \vec{v}}{||\vec{v}||^2} \cdot \vec{v} = 
    \frac{\vec{v} \cdot \vec{u}}{||\vec{u}||^2} \cdot \vec{u} \Rightarrow
    \frac{\vec{u} \cdot \vec{v}}{\vec{v}^2} \cdot \vec{v} = 
    \frac{\vec{u} \cdot \vec{v}}{\vec{u}^2} \cdot \vec{u} 
    \hspace{0.321cm} \overset{*}{\Rightarrow} \hspace{0.321cm}
    \frac{\vec{u} \cdot \vec{u}}{\vec{u} \cdot \vec{u}} \cdot \vec{u} = 
    \frac{\vec{u} \cdot \vec{u}}{\vec{u} \cdot \vec{u}} \cdot \vec{u} \Rightarrow 1=1
\]

\[
    \textbf{\underline{Δεν ξέρουμε αν η πρόταση (α) είναι σωστή} αφού για να ισχύει θα πρέπει:} \hspace{0.321cm} \vec{u} = \vec{v} \hspace{0.321cm} (*)
\]

\[
    \beta) \hspace{0.321cm}
    (\vec{u} \cdot \vec{i})^2 + (\vec{u} \cdot \vec{j})^2 + (\vec{u} \cdot \vec{k})^2 =
    ((x,y,z) \cdot (1,0,0))^2 + ((x,y,z) \cdot (0,1,0))^2 + ((x,y,z) \cdot (0,0,1))^2 =
    x^2 + y^2 + z^2 = ||\vec{u}||^2
\]

\[
    \textbf{Τελικά, \underline{η πρόταση (β) είναι σωστή} ..!}
\]

\[
    \gamma) \hspace{0.321cm} \textbf{Έστω} \hspace{0.321cm}
    \vec{u} = (1,0,0), \vec{v} = (0,1,0), \vec{w} = (1,0,0)
    \hspace{0.321cm} \textbf{τότε} \hspace{0.321cm}
    \vec{u} \cdot \vec{v} = 0, \vec{v} \cdot \vec{w} = 0, \vec{u} \cdot \vec{w} = 1.
\]

\[
    \textbf{Τελικά, \underline{η πρόταση (γ) είναι λανθασμένη} ..!}
\]

\[
    \delta) \hspace{0.321cm} \textbf{Έστω} \hspace{0.321cm}
    \vec{u} = (u_1, u_2, u_3)
    \hspace{0.321cm} \textbf{που είναι κάθετο στο} \hspace{0.321cm}
    <1,1,1>
    \hspace{0.321cm} \textbf{τότε ισχύει} \hspace{0.321cm}
    u_1 \cdot 1 + u_2 \cdot 1 + u_3 \cdot 1 = 0
\]

\[
    \textbf{Όμως, η εξίσωση} \hspace{0.321cm}
    u_1 + u_2 + u_3 = 0 \hspace{0.321cm} \textbf{θυμίζει εξίσωση επιπέδου και όχι ευθείας ...}
\]

\[
    \textbf{Τελικά, \underline{η πρόταση (δ) είναι λανθασμένη} ..!}
\]

\[
    \epsilon) \hspace{0.321cm} \textbf{Εαν} \hspace{0.321cm}
    proj_v(\vec{u}) = 0 \Rightarrow 
    \frac{\vec{u} \cdot \vec{v}}{||\vec{v}||^2} \cdot \vec{v} = 0
    \hspace{0.321cm} \textbf{τότε πρέπει} \hspace{0.321cm}
    \vec{u} \cdot \vec{v} = 0 \hspace{0.321cm} \textbf{ή} \hspace{0.321cm} \vec{v} = 0 \hspace{0.2cm} \textbf{(άτοπο)}
    \hspace{0.2cm} \implies \hspace{0.2cm} \vec{u} \cdot \vec{v} = 0 
\]

\[
    \textbf{Τελικά, αυτό σημαίνει πως είναι κάθετα και άρα \underline{η πρόταση (ε) είναι αληθής} ..!}
\]

\[
    \zeta) \hspace{0.321cm}
    (\vec{u} - \vec{v}) \times (\vec{u} + \vec{v})
    \hspace{0.175cm} \overset{\textbf{επιμεριστική}}{\underset{\textbf{ιδιότητα}}{=}} \hspace{0.175cm}
    \vec{u} \times \vec{u} + \vec{u} \times\vec{v} - \vec{v}\times\vec{u} - \vec{v} \times \vec{v}
    \hspace{0.175cm} \overset{(\vec{u} \cdot \vec{u} = 0)}{\underset{(\vec{v} \cdot \vec{v} = 0)}{=}} \hspace{0.175cm}
    \vec{u} \times \vec{v} - \vec{v} \times \vec{u}
    \hspace{0.175cm} \overset{\textbf{αντισυμμετρική}}{\underset{\textbf{ιδιότητα}}{=}} \hspace{0.175cm} \vec{u} \times \vec{v} - (- \vec{u} \times \vec{v})
    \hspace{0.5cm}
\]

\[
    \textbf{Τελικά,} \hspace{0.321cm} \Rightarrow \hspace{0.1cm} 2(\vec{u} \times \vec{v})
    \hspace{0.321cm} \textbf{και έτσι \underline{η πρόταση (ζ) είναι σωστή} ..!}
\]

\[
    \textbf{στ)} \hspace{0.321cm} \textbf{Λοιπόν, αφησα τελευταίο το ερώτημα (στ) της άσκησης (8) επειδή δεν έχω την παρα-}
\]

\[
    \textbf{-μικρή ιδέα πως να καταλήξω σε οποιοδήποτε συμπέρασμα. Αν θέλετε, αφού οι περισσότερες}
\]

\[
    \textbf{προηγούμενες ήταν σωστές, θα μαντέψω πως αυτή δεν είναι  :)}
\]

\vspace{1.1cm}
\newpage

% ---------------------------------------------------------------------------------------

\newpage

\begin{textblock}{10}(14.401, 15.501)
    \boxed{\textbf{\LaTeX{}}}
\end{textblock}

\begin{center}
    \vspace*{\fill}
    \textbf{\LARGE{ΤΕΛΟΣ}} \\[0.001cm]
    \textbf{\hspace{0.1cm}(ευχαριστώ)}
    \\[0.1cm]
    \vspace*{\fill}
\end{center}

\end{document}

% ---------------------------------------------------------------------------------------

% The End :)